\section*{ВВЕДЕНИЕ}
\addcontentsline{toc}{section}{ВВЕДЕНИЕ}

% TODO (вторая итерация): развернуть до 2--3 страниц с расширенным обоснованием актуальности.

Услуги доступа к сети Интернет за последние два десятилетия превратились из дополнительного удобства в инфраструктурный сервис: бесперебойное соединение требуется как для работы и образования, так и для повседневного потребления цифрового контента. По данным аналитических обзоров российского телекоммуникационного рынка, на территории Российской Федерации действуют сотни локальных интернет-провайдеров (далее~--~ИП), обслуживающих от нескольких сотен до десятков тысяч абонентов каждый. Для таких операторов автоматизация рутинных операций становится критическим условием рентабельности: чем меньше абонентов приходится на одного сотрудника технической и абонентской службы, тем выше относительная стоимость поддержки одного клиента.

Типовой объем рутинной работы оператора-биллинга включает регистрацию нового абонента, привязку его к тарифному плану, выпуск счета на оплату, контроль фактического поступления платежа, приостановку услуг при просрочке и их восстановление после погашения задолженности. Параллельно технический персонал ведет учет сетевого оборудования, контролирует загрузку каналов, реагирует на обращения по горячей линии. При отсутствии единой автоматизированной системы перечисленные операции выполняются вручную в нескольких разрозненных инструментах: банковский интернет-клиент, табличный процессор, веб-интерфейс маршрутизатора, мессенджер для общения с абонентом. Такое разделение приводит к рассогласованию данных, задержкам реакции на инциденты и нерациональному расходованию рабочего времени.

Маршрутизаторы латвийской компании MikroTik занимают значительную долю на рынке сетевого оборудования малых и средних ИП. Операционная система RouterOS предоставляет программный интерфейс, через который сторонние приложения могут управлять простыми очередями (контроль полосы пропускания), правилами брандмауэра, таблицами адресных списков и прочими подсистемами маршрутизатора. Это создает технологические предпосылки для автоматического согласования тарифных планов абонентской базы с сетевыми ограничениями на самом маршрутизаторе.

\emph{Цель настоящей работы}~--~разработка веб-системы для автоматизированного управления клиентами интернет-провайдера, интегрированной с маршрутизаторами MikroTik по протоколу RouterOS API и обеспечивающей снижение трудозатрат на рутинные операции по обслуживанию абонентов. Для достижения поставленной цели необходимо решить \emph{следующие задачи:}
\begin{itemize}
\item провести анализ предметной области и сравнительный обзор существующих решений в классе систем биллинга и управления интернет-провайдерами;
\item сформировать функциональные и нефункциональные требования к разрабатываемой системе;
\item спроектировать многоуровневую клиент-серверную архитектуру и структуру реляционной базы данных;
\item разработать программный интерфейс прикладного уровня в стиле REST;
\item реализовать подсистему интеграции с маршрутизаторами MikroTik для синхронизации тарифных ограничений и управления правилами брандмауэра;
\item реализовать подсистему автоматического биллинга, включающую формирование счетов, приостановку услуг абонентам-должникам и автоматическое восстановление сервиса при погашении задолженности;
\item реализовать подсистему регистрации обращений абонентов сотрудниками технической поддержки и ведения переписки между сотрудниками по обращению;
\item провести модульное и системное тестирование разработанной системы.
\end{itemize}

\emph{Структура и объем работы.} Отчет состоит из введения, 4 разделов основной части, заключения, списка использованных источников и одного приложения. Текст выпускной квалификационной работы равен \formbytotal{lastpage}{страниц}{е}{ам}{ам}.

\emph{Во введении} сформулирована цель работы, поставлены задачи разработки, описана структура работы, приведено краткое содержание каждого из разделов.

\emph{В первом разделе} проводится анализ предметной области: рассматриваются основные бизнес-процессы интернет-провайдера, особенности оборудования MikroTik и протокола RouterOS API, выполняется сравнительный анализ существующих систем биллинга и управления интернет-провайдерами.

\emph{Во втором разделе} формируется техническое задание: определяются цель и назначение разработки, формализуются требования к функциональности и интерфейсу, выполняется моделирование вариантов использования на языке UML, формулируются требования к надежности, информационной безопасности и совместимости.

\emph{В третьем разделе} приведен технический проект: дается общая характеристика проектного решения, обоснован выбор средств и технологий реализации, описано архитектурное построение системы, спроектированы реляционная база данных и программный интерфейс, детализированы подсистемы интеграции с MikroTik и автоматизации биллинга.

\emph{В четвертом разделе} приведен рабочий проект: описаны структура серверной и клиентской частей реализованной системы, перечислены ключевые классы и компоненты, представлены результаты модульного и системного тестирования с приведением контрольных примеров работы.

\emph{В заключении} излагаются основные результаты, полученные в ходе выполнения выпускной квалификационной работы, и формулируются перспективы дальнейшего развития системы.

В приложении~А представлены фрагменты исходного кода разработанной системы.
